\documentclass[10pt,letterpaper]{article}
\usepackage[margin=0.78in,headheight=17pt]{geometry}
\usepackage[T1]{fontenc}
\usepackage[utf8]{inputenc}
\usepackage{lmodern,microtype,parskip}
\usepackage{amsmath,amssymb}
\usepackage{xcolor,booktabs,longtable,tabularx,array,enumitem}
\usepackage[most]{tcolorbox}
\usepackage{fancyhdr,lastpage,titlesec,hyperref,listings}
\usepackage{tikz}
\usetikzlibrary{arrows.meta,positioning,shapes.geometric}
\definecolor{navy}{HTML}{17324D}
\definecolor{bluegray}{HTML}{43657D}
\definecolor{lightblue}{HTML}{EAF2F8}
\definecolor{linegray}{HTML}{CBD5DF}
\definecolor{softgray}{HTML}{F3F5F7}
\definecolor{accent}{HTML}{8E3B46}
\hypersetup{colorlinks=true,linkcolor=navy,urlcolor=bluegray,pdfborder={0 0 0}}
\urlstyle{same}
\titleformat{\section}{\Large\bfseries\color{navy}}{\thesection}{0.6em}{}
\titleformat{\subsection}{\large\bfseries\color{bluegray}}{\thesubsection}{0.5em}{}
\setlist[itemize]{leftmargin=*,itemsep=3pt,topsep=3pt}
\setlist[enumerate]{leftmargin=*,itemsep=3pt,topsep=3pt}
\setlength{\emergencystretch}{3em}
\renewcommand{\arraystretch}{1.22}
\setlength{\tabcolsep}{4pt}
\pagestyle{fancy}\fancyhf{}
\lhead{Secure Programming \& SAST Engineering}
\rhead{Engineering Documentation Suite}
\cfoot{\thepage\ / \pageref{LastPage}}
\newcolumntype{Y}{>{\raggedright\arraybackslash}X}
\newcommand{\checkbox}{\(\square\)}
\newcommand{\blankline}{\noindent\rule{\linewidth}{0.35pt}\par\vspace{4pt}}
\newcommand{\writearea}[1]{\par\begin{tcolorbox}[colback=white,colframe=linegray,boxrule=0.5pt,height=#1,sharp corners]\mbox{}\end{tcolorbox}\par}
\newtcolorbox{focusbox}[1]{colback=lightblue,colframe=bluegray,title=\textbf{#1},fonttitle=\bfseries,breakable}
\newtcolorbox{warnbox}[1]{colback=softgray,colframe=accent,title=\textbf{#1},fonttitle=\bfseries,breakable}
\lstset{basicstyle=\ttfamily\footnotesize,breaklines=true,frame=single,rulecolor=\color{linegray},backgroundcolor=\color{softgray},showstringspaces=false,columns=fullflexible,keepspaces=true,upquote=true}

\begin{document}
\begin{titlepage}
\vspace*{1.0in}
{\color{bluegray}\Large Analysis engineering lab}\\[1.0em]
{\color{navy}\Huge\bfseries SARIF Normalization and Evidence}\\[0.7em]
{\Large Scan-result interchange, fingerprints, disposition and reproducibility}\\[2.0em]
\rule{\linewidth}{1.0pt}
\vfill
\textbf{Source discipline:} Book-derived foundations are identified separately from current engineering extensions.\\
\textbf{Edition:} 29 October 2026
\end{titlepage}
\tableofcontents
\clearpage
\section{Goal}
Interpret, validate and triage SARIF 2.1.0 output from a miniature synthetic scan. SARIF is an interchange format, not a vulnerability confirmation engine.
\section{Inspect the fixture}
Open \texttt{examples/sarif/synthetic-results.sarif}. Locate \texttt{version}, \texttt{runs}, \texttt{tool.driver.rules} and \texttt{results}. For each result identify rule ID, message, relative URI, location, severity and the reviewable code hypothesis. The sample deliberately uses invented source paths and no production secrets.
\section{Tasks}
\begin{enumerate}
\item Validate JSON parsing and the expected SARIF schema version.
\item Explain why \texttt{ruleId} and a physical location alone might be insufficient for stable deduplication across edits.
\item Specify a deterministic fingerprint strategy and categorize results by analyzer and language to avoid collision.
\item Create a human triage record linking finding, immutable commit and reviewer evidence; do not overwrite raw scan output with judgments.
\item State which SARIF fields GitHub may ignore or restrict and verify repository eligibility before real upload.
\end{enumerate}
\section{Gate criteria}
A compliant sample parses; findings are accounted for; one valid false-positive scenario is justified; and the student reports parser success and alert validity as separate claims. Uploading to a production security dashboard is neither required nor appropriate for this fixture.
\end{document}
